\section{模型评价与改进}

\subsection{方法优势与适用条件}

问题一预测品评等级，问题二找出关键化合物，问题三生成风格条件化合物谱并迁移，三个问题各有一条可独立验证的路径。
问题一以逐品评员原始评分的啤酒均分为目标和真值，不依赖附件四从标准化值还原 0 到 6 分的换算；只用附件换算的对照 C1 端到端等级 MAE 为 0.786，主模型为 0.479，差距可直接读出。
问题一的主指标、主输出、成功标准与对照胜出时的处理都在正式运行前写定，正式运行后按声明执行，没有按结果切换模型或规则。
问题二的稳定性用日期组重抽量化，模拟还原实验只用化学矩阵，独立于感官数据，回应了根据结果调整规则的质疑。
问题三的迁移把题面的不恶化要求写成化学层与感官层两层可检验的约束，并在潜空间与原始尺度两处复查。

\subsection{品评信度、关联解释与生成及迁移方法的局限}

化学数据只解释品评小组均分的一小部分方差，连续 $R^2$ 为 0.184，信度（同一批评分随机分两半后相关的外推值）中位数为 0.53，后者给出可解释方差的上限，两者并列说明预测能力接近测量噪声的上限。
问题一的主输出向中间收缩，几乎不给出 5 到 6 级，而真值中这两个等级占 0.968\%。
随机森林对照的等级 MAE 为 0.472，低于主模型的 0.479，但 QWK 为 0.251，低于主模型的 0.290；主模型是运行前选定的模型，不是外层比较中误差最低的模型。
问题二的第 3 位基本不稳定，入选频率按名次约为 0.60、0.36、0.22，与主要对照方法 M3 的 Top-1 一致比例只有 0.45，这是不同估计对象的差异，不是某一个错了。
问题三的生成样本比真实啤酒更接近风格均值，风格判别器对真实啤酒的识别准确率为 0.550，对生成样本为 0.481，生成样本围绕收缩后的风格均值分布。
问题三的感官层约束依赖问题一在缺陷维度上的预测，这些维度多为零分主导、信度低，化学层约束才是主要保证，感官层是模型层面的检查。

\subsection{数据补充、模型改进与实验验证方向}

问题一的模型族选择可把梯度提升树与随机森林连同扩展后的网格一起纳入内层选择，重新做嵌套评价，这需要问题一主线以及依赖其选参的问题二主要对照与问题三一起重算。
高等级尾部可改用有序回归或分位数回归直接校准，改善 5 到 6 级的召回。
问题三的潜变量相关可改用计入截断的秩估计（截断连续变量的 Kendall $\tau$ 桥函数，Fan 等人 2017；Yoon 等人 2020），使下限状态的共现可被生成，并按预测分布抽样、计入风格均值的估计误差；这两项改变主模型，须纳入内层选参后重新嵌套评价。
